% Part: first-order-logic
% Chapter: model-theory
% Section: overspill

\documentclass[../../../include/open-logic-section]{subfiles}

\begin{document}

\olfileid{mod}{bas}{ove}
\olsection{حد توں اگے پھیلاؤ (Overspill)}

\begin{thm}
\ollabel{overspill} جے جملیاں دے سیٹ $\Gamma$ دے من مانے وڈے
محدود نمونے ہون، تاں اوہدا اک لامتناہی نمونہ وی ہوندا اے۔
\end{thm}

\begin{proof}
$\Gamma$ دی زبان وچ قابلِ شمار لامتناہی تعداد دے نویں مستقل
$c_0$، $c_1$، \dots شامل کرو تے سیٹ $\Gamma \cup \{c_i \neq c_j :
i \neq j\}$ ویکھو۔ ایہ کہنا کہ $\Gamma$ دے من مانے وڈے محدود نمونے
نیں، ایس دا مطلب اے کہ ہر $m >0$ لئی کوئی $n\ge m$ اے جیہدے لئی
$\Gamma$ دا ایداں دا نمونہ اے جس دے دائرے وچ $n$ عنصر نیں۔ ودھے سیٹ دا
کوئی وی محدود حصہ چنو۔ اوہدے وچ صرف محدود تعداد دے نویں مستقل آندے
نیں؛ اصل سیٹ دا اِنّے عنصر رکھن والا یا اوہدوں وڈا محدود نمونہ چن کے
اوہناں مستقل نشانیاں نوں وکھ وکھ عنصر دیو۔ باقی نویں مستقل اک من مرضی
دے عنصر نوں نام دے سکدے نیں؛ اصل زبان دی تشریح نہیں بدلنی۔ ایس لئی
چنے محدود حصے دے اصل جملے تے عدمِ برابریاں اکٹھیاں سچ ہون گیاں۔ ایس توں
ایہ نکلدا اے کہ $\Gamma \cup \{c_i \neq
c_j : i \neq j\}$ محدود طور تے پورا ہون جوگا اے۔ محدود گواہی دے
قضیے مطابق $\Gamma
\cup \{c_i \neq c_j : i \neq j\}$ دا اک نمونہ $\Struct M$
اے جس دا دائرہ لازماً لامتناہی اے، کیوں جے ایہ ساریاں عدمِ برابریاں
$c_i \neq c_j$ پوری کردا اے۔
آخر وچ نویں مستقل نشانیاں دی تشریح بھلا کے اصل زبان والی ساخت لے لو؛
دائرہ اوہی لامتناہی رہندا اے تے اصل سیٹ دے سارے جملے وی سچے رہندے نیں۔
\end{proof}

\begin{prop}
\ollabel{inf-not-fo}
کسی وی پہلے درجے دی زبان دا کوئی جملہ $!A$ ایسا نہیں کہ اوہ
!!a{structure}~$\Struct M$ وچ اُدوں تے صرف اُدوں سچ ہووے جدوں
اوس !!{structure} دا دائرہ $\Domain{M}$ لامتناہی ہووے۔
ایتھے اکّو جملے نال ساریاں تے صرف لامتناہی ساختاں دی پہچان دی گل اے؛
جملیاں دے پورے لامتناہی سیٹ نال ایہ شرط بیان کر سکّن دی نفی نہیں۔
\end{prop}

\begin{proof}
جے ایداں دا کوئی $!A$ ہوندا، تاں اوہدی نفی $\lnot !A$ ہر تے
صرف محدود !!{structure}s وچ سچی ہوندی۔ ایس لئی اوہدے من مانے وڈے
محدود نمونے ہوندے، پر کوئی لامتناہی نمونہ نہ ہوندا؛ ایہ
\olref{overspill} دے خلاف اے۔
\end{proof}

\end{document}
